\documentclass[11pt]{article}
\usepackage[T1]{fontenc}
\usepackage[margin=1in]{geometry}
\usepackage{amsmath,amssymb,amsthm,mathtools,microtype}
\usepackage[colorlinks=true,allcolors=blue]{hyperref}
\newtheorem{theorem}{Theorem}[section]
\newtheorem{lemma}[theorem]{Lemma}
\newcommand{\mc}{\mathsf{MCSP}}
\newcommand{\ad}{\widetilde{\deg}}
\newcommand{\E}{\mathbb{E}}
\newcommand{\F}{\mathbb F_2}
\title{The approximate degree of Boolean minimum circuit size}
\author{Research manuscript}\date{22 November 2026}
\begin{document}\maketitle
\begin{abstract}
How well can a real polynomial approximate the test that an $N$-entry truth table has a small Boolean circuit? Write $n=\log_2 N$ and $T=s\log_2 s$ for a gate budget $s$. We obtain matching bounds of order $\sqrt{NT/(1+\log_2(N/T))}$ throughout $(2+\eta)n\le s\le(1-\epsilon)N/n$ for fixed $\eta,\epsilon>0$. This is also the bounded-error quantum query complexity of the same classical decision problem. We also determine the dependence on the worst-case error probability. The lower bound comes from a polynomial in one complex variable. We sample a sparse mask in each row and repeat it after shifting its entries by a secret label. In the easy experiment many row labels share a short description. Independent labels make the table hard. Matching the low-order marginals of the labels forces the polynomial to vanish to high order at the origin. The acceptance gap at a nearby positive point then forces a large degree.
\end{abstract}

\section{The decision problem}
A truth table $x\in\{0,1\}^N$ specifies a Boolean function on $n$ address bits. Write $C(x)$ for its minimum circuit size with binary AND and OR, unary NOT, free constants and unrestricted fanout. The fixed-budget decision function is $\mc_{n,s}(x)=\mathbf1[C(x)\le s]$. Its approximate degree $\ad$ is the least degree of a real polynomial within $1/3$ at every Boolean input. Multilinearization preserves all Boolean evaluations and can only decrease degree. Write $Q$ for bounded-error quantum query complexity with coherent access to the classical input bits. Logarithms are base two unless specified otherwise.
\begin{theorem}\label{th:main}
For fixed $\epsilon,\eta>0$, uniformly throughout $(2+\eta)n\le s\le(1-\epsilon)N/n$, put $T=s\log s$. Then
\[
 \ad(\mc_{n,s})=\Theta_{\eta,\epsilon}\!\left(\sqrt{\frac{NT}{1+\log(N/T)}}\right),
 \qquad Q(\mc_{n,s})=\Theta_{\eta,\epsilon}\!\left(\sqrt{\frac{NT}{1+\log(N/T)}}\right).
\]
The same bounds hold if the inner circuit may use all binary Boolean gates.
\end{theorem}
For $s=N^{a+o(1)}$ with $0<a<1$, the formula has order $\sqrt{Ns}$, with constants depending on $a$. For $n=o(s)$ polynomially bounded in $n$, its order is $\sqrt{Ns\log s/n}$. The usual sparse-family identification algorithm proves the upper bound. For the lower bound, Section~2 gives a general observation about sparse masks. Sections~3 and~4 construct the Boolean truth tables to which it applies.

\paragraph{Counting and the upper bound.}
Let $A_s=\{x:C(x)\le s\}$. Traverse a rooted single-output circuit of $j$ useful gates backwards in depth-first order. Of at most $2j$ predecessor slots, $j-1$ introduce a new gate. The remaining at most $j+1$ pointers need $\log(n+j+2)+O(1)$ bits each. The ordered traversal, slot designations and labels require $O(j)$ more bits. Summing over smaller circuits yields
\[
                \log |A_s|\le(s+1)\log(n+s+2)+O(s+\log n).               \tag{1}
\]
This is $(1+o(1))T$ in the theorem's range. Kothari's algorithm \cite{koth} identifies a promised member of a family of $M$ possible $N$-bit inputs using $O(\sqrt{N\log M/(1+\log(N/\log M))})$ queries when $N\le M\le 2^N$. The $M<N$ case costs $O(\sqrt N)$. Run the capped algorithm and verify its proposed table with a search for a mismatched entry. This costs $O(\sqrt N)$ additional queries and gives a bounded-error decision algorithm even off the identification promise. If $T$ is a constant fraction of $N$, reading the entire table also suffices. A $q$-query algorithm has an acceptance probability polynomial of degree at most $2q$ \cite{beals}. Thus one lower bound for ordinary approximate degree suffices for both claims.

\section{Random masks and a polynomial in one variable}
Factor $N=qkH$ into powers of two. Index a table by row $i\in[k]$, copy $c\in\{0,\ldots,q-1\}$ and location $j\in\F^{h}$ where $H=2^h$. Independently in each row choose a random mask $B_i$ by including each of its $H$ locations with probability $t$. Put $Y_{i,0}=0$. Other shifts $Y_{i,c}$ are independent of the masks. The input table is
\[
             X(i,c,j)=\mathbf1[j\mathbin\oplus Y_{i,c}\in B_i].              \tag{2}
\]
At a row with a singleton mask, the first copy identifies the one marked location. Each other copy identifies the translated location and hence its shift.

\begin{lemma}\label{lem:analytic}
Compare two laws of the row vectors $(Y_{i,1},\ldots,Y_{i,q-1})$ that agree on every set of at most $D$ rows. Suppose the acceptance probabilities of $\mc_{n,s}$ on (2) at $t_0=\lambda/H$ are $1-o(1)$ and $o(1)$ respectively, where $\lambda$ stays in a compact positive interval. If $H\to\infty$ and $D\to\infty$, then $\ad(\mc_{n,s})=\Omega(D\sqrt H)$.
\end{lemma}
\begin{proof}
For a global degree-$d$ approximator $P$, let $R(t)$ be the difference between its expectations in the two experiments. For every $0\le t\le1$ these are true distributions on Boolean tables, so $|R(t)|\le5/3$. At $t_0$ the gap is at least $1/3-o(1)$.

Fix the labels and a multilinear monomial of $P$. Its expectation over masks is $t^u$, with $u$ the number of distinct underlying pairs $(i,j\oplus Y_{i,c})$ at which the monomial requires a mark. In particular $u$ is at least the number of distinct rows touched and at most the monomial degree. A monomial touching at most $D$ rows has identical expectations under the two label laws. Every other monomial contributes only powers above $D$. Thus $R$ has degree at most $d$ and a zero of order at least $D$ at the origin.

For completeness, if a polynomial is bounded by $M_0$ on $[0,1]$, its maximum on $|z|\le r$ for small $r>0$ is at most $M_0\exp(Cd\sqrt r)$ for an absolute $C$. Indeed transform the interval to $[-1,1]$ by $w=2z-1$ and use the inverse Joukowski map $w=(u+u^{-1})/2$ with $|u|\ge1$. Applying the maximum principle to the transformed polynomial divided by $u^d$ outside the unit circle bounds it by $M_0|u|^d$. For $|w+1|\le2r$, the formula $u=w\pm\sqrt{w^2-1}$ gives $|u|\le1+C\sqrt r$. Apply the maximum principle on the disk to $R(z)/z^D$, taking $r=4t_0$. Then
\[
            |R(t_0)|\le (5/3)4^{-D}\exp(C'd\sqrt{t_0}).
\]
Its positive lower bound proves the lemma. An alternative proof follows from Vladimir Markov's higher derivative inequality. Closely related polynomial growth estimates are used in search direct-product bounds \cite{ksw}.
\end{proof}

\section{Superlinear gate budgets}
Assume first that $n=o(s)$ and $T=o(N)$. Let $k$ be the smallest power of two on the increasing branch for which $kh\ge T$, where $h=\log(N/(2k))$. Set $q=2$. Then $T\le kh\le2T$, $h\to\infty$, and $k=\Theta(T/(1+\log(N/T)))$. To check the last relation, write $v=\log(N/T)\to\infty$. At $k$ of order $T/v$, $h=v+\log v+O(1)$. Also $k\ge T/n\gg\log n$. Take $\lambda=0.8T/(kh)$, so $0.4\le\lambda\le0.8$.

\paragraph{The private mask.}
The mask table $B$ takes $\log(kH)=n-1$ inputs. At $t_0$ its weight $L$ is $(\lambda+o(1))k$ with probability tending to one. Lupanov's sparse synthesis theorem \cite{lup,kors} gives, for every such table, a circuit of size $(1+o(1))E/\log E$ where $E=\log\binom{kH}{L}$. Here $L/\log n\to\infty$, $E=(\lambda+o(1))kh=(0.8+o(1))T$, and $\log E\sim\log s$. The cost is $(0.8+o(1))s$. Chashkin's synthesis bounds on prescribed sets also apply \cite{chash}.

\paragraph{Cheap shifts with matching marginals.}
For any sufficiently small fixed $\rho>0$ we construct full $h$-bit shifts that are jointly uniform on any $D=\lfloor\delta k\rfloor$ rows, $\delta=\delta(\rho)>0$. Their evaluation circuits have size $O(\rho s)+o_\rho(s)$ with an absolute constant in the first term.

If $\log k\ge (\log s)/3$, start with a binary linear space of words of length $k$ whose dual has distance at least $\delta k$ and each of whose words costs at most $(9\rho+o_\rho(1))k/\log k$ gates. Here is a construction. Factor $k=b_1 b_2$ with the two powers of two within a factor of two. For $a=1,2$, choose a $\lceil\rho b_a\rceil$ by $b_a$ parity-check matrix with kernel distance at least $\Delta b_a$ for small fixed $\Delta>0$; a union bound gives such a matrix. Write its columns as $g_{a,j}$. The code consists of arrays
\[
        X_{ij}=\langle U_i,g_{2,j}\rangle\oplus\langle V_j,g_{1,i}\rangle
\]
with independently chosen vectors $U_i,V_j$. In its dual, every row and column belongs to the respective kernel. Any nonzero word therefore has weight at least $\Delta^2 k$. The four ordinary lookup tables for the vectors and columns each cost at most $(1+o(1))b_a/\log b_a$ per bit by Lupanov synthesis. Their total is at most $(2+o(1))\rho (b_1+b_2)^2/\log k\le(9\rho+o_\rho(1))k/\log k$. The dot products cost $O(b_1+b_2)$ additional gates, absorbed in the error. Take $h$ independent words from this space. Their total cost meets the claim since $kh\le2T$.

If $\log k<(\log s)/3$, use a binary matrix with $k$ columns and $m=\lceil\rho k\rceil$ rows in which each column has four ones and every set of at most $\delta k$ columns is independent. To construct one, choose four random distinct rows in each column independently. A set of $u$ columns without a unique neighbor uses at most $2u$ rows, an event with union bound $\binom{k}{u}\binom{m}{2u}(2u/m)^{4u}\le(C_\rho u/k)^u$ when $2u\le m$. Sum up to $\sqrt k$ and then to sufficiently small $\delta k$. Define each shift as the XOR of the four indexed words from $m$ independent uniform $h$-bit words.

To evaluate the shifts, compute the row-address indicators in $O(k\log k)$ gates. For each of four neighbor positions group those indicators by the selected seed, forming a one-hot list of length $m$ in $O(k)$ gates. Split the list into groups of $\lfloor\tfrac12\log h\rfloor$ entries. Generate every subset OR in each group, then form each output bit as the OR of one subset result per group. The four lookups together cost $O(\rho kh/\log h+h+k\log k)$. The case assumption implies $h\ge s^{2/3}$ for large $s$, while $h\le n=o(s)$ and $k\log k=o(s)$. Choosing $\rho$ sufficiently small proves the claim. Computing the selected shifted address and using the mask then takes only $O(h)$ extra gates. All but an $o(1)$ fraction of easy-law outcomes lie in $A_s$.

\paragraph{The hard reference.}
Choose the $k$ shifts independently and uniformly. Besides the typical mask weight, with probability $1-o(1)$ the number $J$ of singleton rows is $(\lambda e^{-\lambda}+o(1))k$, by independent rows and $H\to\infty$. Any fixed typical output identifies its first-copy mask. Its point probability is at most
\[
      (\lambda/H)^L(1-\lambda/H)^{kH-L} H^{-J}
      \ \le\ 2^{-[\lambda(1+e^{-\lambda})+o(1)]kh}.
\]
Indeed the singleton rows identify independent shifts. Since $0.8(1+e^{-0.8})>1$, the exponent exceeds $T$ by a fixed fraction. Summing over (1) proves that the probability of an easy output tends to zero. In both experiments the interpolation in Lemma~\ref{lem:analytic} uses the original, unconditioned distributions. Applying it gives $\Omega(k\sqrt H)=\Omega(\sqrt{Nk})$.

\section{The two remaining ranges}
\paragraph{A fixed linear gate ratio.}
Suppose $(2+\eta)n\le s\le Cn$ with $C$ fixed. Choose $k$ a power of two between $\log n$ and $2\log n$. Choose a constant power of two $q$ sufficiently large in terms of $C$, put $H=N/(qk)$, and set $\lambda=(\log n)/(3k)$. Then $1/6\le\lambda\le1/3$ and $h=n-O(\log\log n)$. For each $c=1,\ldots,q-1$ generate its shifts across rows by the four-neighbor code of Section~3, independently for different copies. Choose its fixed rate $\rho$ small enough that $(q-1)\lceil\rho k\rceil\le (\log n)/8$ for large $n$. Every set of at most $\delta k$ complete row vectors is jointly uniform.

With high probability the first-copy masks contain $L\le (\log n)/2$ marks in total. The suffix columns of all $qL$ positive addresses have at most $2^{L+(q-1)\lceil\rho k\rceil}$ distinct binary patterns. Indeed, their entries in one suffix column are determined by the entries for $L$ private marks and the seed bits in that same column for all shifts. Group the $h$ raw suffix bits by identical pattern. In each group of size $b$, test both all-zero and all-one using at most $2b$ AND/OR/NOT gates. Each positive address can then be recognized by testing its row, its copy, and the prescribed value in each group. OR the positive-address tests. The total cost is at most
\[
            2h+O\!\left(qL(2^{L+(q-1)\lceil\rho k\rceil}+\log(qk))\right)
              \le 2n+o(n).
\]
This fits the budget. In the independent reference each singleton row reveals $q-1$ independent shifts. The negative logarithm of a typical table's point probability is at least
$[1+(q-1)e^{-\lambda}+o(1)]\lambda k h$.
Choose $q$ so $(1+(q-1)e^{-1/3})/3>C+1$. The circuit census (1) makes the reference hard with high probability. Lemma~\ref{lem:analytic} now gives $\Omega(\sqrt{N\log n})$.

\paragraph{A fixed positive fraction of the Shannon budget.}
If $T/N$ tends to a positive constant and $s\le(1-\epsilon)N/n$, use the binary code of Section~3 directly on $N$ table entries, with sufficiently small constant $\rho$ so every codeword has cost at most $s$. Its uniform law and the uniform law on all $N$ bits have the same joint marginals through order $\delta N$. Every polynomial of degree at most $\delta N$ has identical expectations. Yet a random table has size exceeding $s$ with probability $1-o(1)$ by (1). This gives $\Omega(N)$ directly. For the uniform statement of Theorem~\ref{th:main}, any sequence allegedly violating a positive lower constant has a subsequence with $s/n$ bounded or tending to infinity and $T/N$ convergent. One of the three lower-bound arguments applies to that subsequence with a positive constant, a contradiction.

\section{Small error on the full truth table}\label{sec:confidence}
Write $Q_\delta$ for quantum queries with error at most $\delta$ and $\deg_\delta$ for the ordinary approximate degree at this tolerance. A small error can be obtained at a lower cost than independent repetition of the entire decision algorithm.

\begin{lemma}[Identifying a word with small error]\label{lem:small-error-oip}
Let $A$ be a nonempty known family of $N$-bit words. Put $B=\min\{N,1+\log|A|\}$, $K=1+B/(1+\log(N/B))$ and $\ell=\log(1/\delta)$ for $0<\delta\le1/3$. There is a quantum algorithm with
\[
                O\bigl(\min\{N,\sqrt{N(K+\ell)}\}\bigr)
\]
queries which returns the input word when it belongs to $A$, or rejects when it lies outside $A$, with probability at least $1-\delta$ on every input. Computation on the known family is free.
\end{lemma}
\begin{proof}
Set $r=\lceil K+\ell\rceil$. Read the whole word if $r\ge N/16$. Otherwise put $a=\sqrt{r/N}$, $\gamma=c a\log(2/a)$, where $c>0$ is a sufficiently small absolute constant. Keep the words consistent with each explicitly read bit. Reject when no survivor remains. Whenever a coordinate has minority frequency at least $\gamma$, query it. Each such query retains at most a $(1-\gamma)$ fraction, so at most $O(B/\gamma)$ occur.

When all coordinates have smaller minority frequency, store the survivors' coordinatewise majority $m$ as one vote. Search for a disagreement with $m$, verifying any returned index. A search using $O(1/a)$ queries finds a disagreement with probability at least $3/4$ whenever $d(x,m)>r/8$. This is bounded-error quantum search with a known lower bound on the marked fraction \cite{dw}. A verified disagreement retains at most a $\gamma$ fraction. There are at most $J=1+\lceil B/\log(1/\gamma)\rceil$ successful searches, including a last search that empties the family. A singleton keeps voting by the same rule.

Stop at $q=\lceil C(J+\ell)\rceil$ votes for a sufficiently large absolute $C$. A vote is bad when $d(x,m)>r/8$. At each bad vote, conditional on the entire history, search succeeds with probability at least $3/4$. Couple these success indicators with independent Bernoulli$(3/4)$ trials indexed by the successive bad votes, padding unused trials after stopping. A run reaching $q$ votes with at least $u=\lceil q/4\rceil$ bad votes has at most $J$ successes among the first $u$ trials. The binomial tail is at most $\delta$. Otherwise the coordinatewise majority $z$ of the stored words differs from $x$ in at most $r/2$ positions. Every wrong coordinate accounts for at least $q/4$ disagreements among good votes, whereas all good votes together account for at most $qr/8$ disagreements.

To correct $z$, run exact quantum search successively for $j=r,r-1,\ldots,1$ marked positions among the as yet uncorrected discrepancies. Each invocation costs $O(\sqrt{N/j})$ and any reported index is verified. If the remaining number is $j$ the search succeeds with certainty. If it is smaller, verification can only reduce it. Induction thus corrects every discrepancy whenever initially at most $r$ were present \cite{dw}. Check offline whether the result belongs to $A$.

The cost is $O(B/\gamma+(J+\ell)/a+\sqrt{Nr})$. Here $\log(1/\gamma)=\Theta(\log(2/a))$ and $B=O(r\log(2/a))$. For the latter inequality put $u=B/(1+\log(N/B))\le B$. Then $u(1+\log(N/u))\ge B$, and $v(1+\log(N/v))$ is increasing on the interval concerned, $u\le v\le N/16$. Substitution gives the stated bound.
\end{proof}

\begin{theorem}[The full error dependence]\label{th:small-error}
In the gate range of Theorem~\ref{th:main}, put $K_N=T/(1+\log(N/T))$. Uniformly for $0<\delta\le1/3$,
\[
 \deg_\delta(\mc_{n,s}),\ Q_\delta(\mc_{n,s})
      =\Theta_{\epsilon,\eta}\!\bigl(\min\{N,\sqrt{N(K_N+\log(1/\delta))}\}\bigr).
\]
\end{theorem}
\begin{proof}
The circuit census (1) and Lemma~\ref{lem:small-error-oip} give the algorithm. Its acceptance polynomial has at most twice its query count as degree. Theorem~\ref{th:main} gives the first lower-bound term. For the second, start with $A_s$ and repeatedly query a coordinate on which the survivors disagree, retaining the smaller nonempty side. Within $\lfloor\log|A_s|\rfloor$ steps just one word $v$ remains. Fix these coordinates to their values in $v$ and complement each remaining coordinate on which $v$ is one. The restricted membership function is NOR on $N-\lfloor\log|A_s|\rfloor$ or more inputs. The census and the fixed Shannon margin leave $\Omega_\epsilon(N)$ such inputs. NOR on $d$ bits has $\delta$-approximate degree $\Theta(\min\{d,\sqrt{d\log(1/\delta)}\})$ \cite{dw}. Restrictions preserve the degree upper bound, proving the claim for either inner gate basis.
\end{proof}

\section{Testing at a prescribed Hamming distance}\label{sec:test}
A tester must accept $x\in A_s$ and reject tables with $d(x,A_s)\ge\xi N$, where $d$ counts differing entries. Its output elsewhere is unrestricted. Let $R_\xi,Q_\xi$ be the classical randomized and quantum query costs with pointwise error at most $1/3$ on these two cases. Computation between queries is free.
\begin{theorem}\label{th:testing}
Fix $\eta,\epsilon>0$. For some $\xi_0>0$ and every $1/N\le\xi\le\xi_0$, in the gate range of Theorem~\ref{th:main}, write $L_\xi=1+\log(1/\xi)$ and $K_N=T/(1+\log(N/T))$. Then
\begin{align*}
 R_\xi &=\Theta_{\eta,\epsilon}\!\left(\min\{N,\, T/(\xi L_\xi)\}\right),\\
 Q_\xi &=\Theta_{\eta,\epsilon}\!\left(\min\{\sqrt{NK_N},\,T/(\sqrt\xi L_\xi)\}\right).
\end{align*}
\end{theorem}
Thus at constant proximity in this range both models use order $s\log s$ queries. As the allowed distance decreases, knowing the input is far on a no-instance still helps until $\xi$ reaches order $K_N/N$.

\paragraph{An upper bound for any finite family.}
Let $A$ be a known class of at least two binary strings of length $N$ and let $B=\log |A|$. Choose $a=\xi$ classically and $a=\sqrt\xi$ quantumly. Searching for a coordinate at which two tables differ can be done with a constant probability of finding a mismatch using $O(1/a)$ queries whenever they disagree on at least $\xi N/3$ coordinates. Use random sampling or quantum search with the given lower bound on the marked fraction, always verifying any proposed coordinate. Set $\gamma=c a\log(2/a)$ with a sufficiently small absolute $c$ so $\gamma<1/2$.

Keep only candidates consistent with actual queried values. If a coordinate splits the survivors with minority proportion at least $\gamma$, query it. At most $O(B/\gamma)$ such queries occur since each retains at most a $(1-\gamma)$ fraction. Otherwise let $m$ be the coordinatewise majority of the survivors. Cast one vote for acceptance if some survivor lies within $\xi N/3$ of $m$ and for rejection otherwise. Then make one search attempt for a disagreement of $x$ with $m$. A found mismatch retains at most the $\gamma$ fraction. If an attempt finds nothing, the survivors stay in place. When they become empty reject; for a singleton verify its member by a single amplified promise search, distinguishing equality from distance $\xi N$.

Here is a uniform bound on the accumulated error. On a promised input, whenever the vote is wrong, $d(x,m)\ge\xi N/3$. Indeed if $x\in A$ it is still a survivor; if $x$ is far then the triangle inequality applies to the proposed nearby survivor. A fixed survivor set with a wrong vote therefore contributes at most a constant expected number of wrong votes before a mismatch is found. There are at most $1+B/\log(1/\gamma)$ distinct sets at which votes are cast: every successful minority search reduces the survivor count by a factor $\gamma$. Stop after a large constant times $1+B/\log(1/\gamma)$ votes or earlier termination. If the vote cap is reached, output the majority. Markov's inequality bounds the probability of a wrong majority, with any desired fixed error margin by choice of cap. Balanced queries and witnessed mismatches never remove the target on a yes-input. This proves the query bound
\[
 O\!\left(\frac1a+\frac{B}{a\log(2/a)}\right).
\]
Use (1), the exact upper bounds, and $T\gg\log N$ to obtain both claimed upper bounds.

\paragraph{Separation from a Hamming neighborhood.}
For an ambient length $M=2^m$ under Theorem~\ref{th:main}'s gate hypotheses, set $K_M=T/(1+\log(M/T))$. There is an $\alpha>0$, depending only on the fixed endpoint margins, such that the two experiments in our proof separate $A_s$ from distance at least $\alpha K_M$ using lower bounds $\Omega(M)$ classical queries and $\Omega(\sqrt{M K_M})$ quantum queries. For the distance assertion, enlarge the easy family in the hard-law union bound to its Hamming neighborhood. In the superlinear sparse range a radius $\alpha K_M$ adds at most $(O(\alpha)+o(1))T$ bits to the log count. This follows from $\sum_{j\le r}\binom Mj\le (eM/r)^r$ for $1\le r\le M/2$, and $\log(M/K_M)\sim1+\log(M/T)$ when $T=o(M)$. The hard point-probability exponent has a fixed fractional margin over $T$. The same reasoning applies at bounded linear gate ratios, taking the number of copies large enough that the exponent in Section~4 exceeds $2T$. At a positive Shannon fraction the uniform hard law gives the neighborhood probability bound $|A_s|2^{-M}\sum_{j\le\alpha M}\binom Mj=o(1)$ if binary entropy $H_2(\alpha)<\epsilon/2$. The constants can be taken uniformly by the subsequence argument of Section~4.

For the classical query lower bound, fix a deterministic adaptive $u$-query algorithm and fix the shifts. Each query to (2) either reads a known mask bit or reveals a fresh independent Bernoulli bit. Its expected output at complex parameter $z$ is a polynomial of degree at most $u$ and modulus bounded by $(|z|+|1-z|)^u$. To see this, write it as a sum of transcript weights, retaining the value of a bit after its first read, and sum their absolute values down the binary decision tree. Average over algorithm coins and over the two label laws. Their expectation difference $R$ has $|R(z)|\le2\exp(2u r)$ on $|z|\le r$. The decision polynomial is multilinear of degree at most $u$, so the same touched-row argument as in Lemma~\ref{lem:analytic} gives a zero of order $D$ at zero. At $r=4t_0$, $|R(t_0)|\le 2\cdot4^{-D}\exp(8u t_0)$, while the promise forces a constant gap in the real acceptance probabilities. Thus $u=\Omega(DH)=\Omega(M)$ in the mask experiments. In the dense code experiment, bounded independence already fools every deterministic decision tree of depth at most $\delta M$, and hence their randomized mixtures. For the quantum lower bound in all cases, the decision algorithm gives an acceptance polynomial bounded throughout the Boolean cube, even on inputs outside its promise.

\paragraph{Adapting the domain length.}
Fix $c_0>0$ less than the neighborhood constant $\alpha$, using a Shannon margin of $\min(\epsilon,1/2)$ in that constant. Choose $\xi_0$ sufficiently small in terms of $c_0$. Let $M$ be the greatest power of two at most
\[
                 \min\{N,\ c_0 T/(\xi L_\xi)\}.
\]
If $M<N$, we can ensure $M\ge4T$ and $1+\log(M/T)=\Theta(L_\xi)$ by shrinking $\xi_0$. Since $m/2^m$ is decreasing for $m\ge2$, $s m/M\le1/2$ for large $s$, by checking near $M=4T$. If $M=N$, the original gate margins apply. In both cases $1+\log(M/T)\le L_\xi$, whence $\xi M\le\alpha K_M$. Duplicate a table on $m$ address bits across all assignments to the remaining $n-m$ bits. Its minimum circuit size is unchanged by duplication since any ambient circuit can be restricted. Each such restriction of a size-$s$ circuit must disagree with a hard $m$-variable table on at least $\alpha K_M$ positions. Summing these discrepancies over the slices shows the padded table is $\xi N$-far. One query to it costs one query to the original table, coherently as well as classically. For $M=N$ this recovers the exact lower scale. For $M<N$, substitute $M=\Theta(T/(\xi L_\xi))$ and $K_M=\Theta(T/L_\xi)$ into the preceding separation, completing the proof.

The testing-by-implicit-learning literature gives classical constant-distance upper bounds via finite families on relevant inputs. Alon, Blais, Chakraborty, Garc\'ia-Soriano and Matsliah prove a polynomial-in-$s$ query lower bound for testing the entire class of size-$s$ circuits at constant distance (their Section~7.4, Corollary~2.2). Their work also studies function isomorphism \cite{iso,extractors}. Our calculation gives both query models at varying proximity for fixed binary fan-in and the gate range in Theorem~\ref{th:testing}. It makes no efficiency assumption on enumerating candidate circuits.

\section{The cost of choosing all queries in advance}
A nonadaptive classical tester chooses its full list before seeing the answers. In the quantum version, it prepares any entangled state, applies the bit oracle to all its query registers in one batch, and performs a measurement. Let $R^{\rm na}_\xi,Q^{\rm na}_\xi$ denote their bounded-error query costs.
\begin{theorem}\label{th:nonadaptive}
For fixed $\eta,\epsilon>0$, let $(1+\eta)n\le s\le(1-\epsilon)N/n$. For every $1/N\le\xi\le\xi_0$, where $\xi_0>0$ is a sufficiently small fixed constant depending on the margins,
\[
              Q^{\rm na}_\xi=\Theta_{\eta,\epsilon}(\min\{N,T/\xi\})
                   =R^{\rm na}_\xi .
\]
Both bounds also hold when the tester must accept every table within $\theta\xi N$ of $A_s$, for any fixed $0\le\theta<1$. The upper-bound constant may depend on $\theta$.
\end{theorem}
\begin{proof}
For the upper bound, query all entries or draw $O_\theta(T/\xi)$ uniform random entries. Accept if some member of $A_s$ has observed disagreement below $(1+\theta)\xi/2$. If a close member exists, its empirical disagreement obeys a multiplicative Chernoff upper-tail bound, using stochastic domination by sampling from a population with disagreement $\theta\xi$. On a far input every member has actual disagreement at least $\xi$. Its lower-tail failure has probability $\exp(-\Omega_\theta(\xi q))$ after $q$ samples, and (1) allows a union bound. Here observed disagreement denotes the sample fraction. Coincidences among sampled entries can be queried only once.

The lower bound uses a randomly chosen block of addresses. First suppose a power of two $M\in[4T,8T)$ fits inside $N$. Even if $M$ lies slightly above $N/2$, this argument applies. Use a uniform word in the binary linear space of Section~3 with length $M$ and small constant rate parameter $\rho$. For some $d=\lfloor\delta M\rfloor$ it is $d$-wise uniform. Every such word, interpreted as a Boolean table on $\log M$ address bits, is computable with at most $(9\rho+o(1))M/\log M$ gates. Take $\rho$ sufficiently small in terms of $\eta$, leaving at least $n+O(1)$ gates. In the hard experiment take $M$ independent fair bits. A fixed small $\alpha>0$ with $H_2(\alpha)<1/4$ ensures this word is at distance at least $\alpha M$ from every $s$-gate table on $\log M$ bits with probability $1-o(1)$. Indeed the circuit census on these bits has logarithm $(1+o(1))T\le(1+o(1))M/4$.

Let $p$ be the smallest reciprocal power of two at least $\max(M/N,\xi/\alpha)$ and choose $\xi_0\le\alpha/2$. Partition the table by a prefix of $\log(1/p)$ address bits. Choose one prefix uniformly. On its block write the chosen word on the last $\log M$ bits, ignoring any intervening bits. Set all other addresses to zero. The selector of a prefix of length $b$ uses at most $b+O(1)$ binary AND, OR and NOT gates: collect all required zeros with one OR followed by NOT, all required ones with one AND, then combine. The final gate combines this selector with the word, so the easy experiment stays within budget. In the hard experiment fix the chosen prefix and each assignment of the intervening bits. Every size-$s$ approximant on the full domain restricts to a size-$s$ approximant on this slice. Summation of its discrepancies gives distance at least $p\alpha N\ge\xi N$. The choice of $p$ gives $M/p=\Theta(\min(N,T/\xi))$.

Consider any one-batch $q$-query quantum tester, purifying its state $|\psi\rangle$ if necessary. For a fixed chosen block, let $\Pi$ keep the computational-basis query tuples with at most $\lfloor d/2\rfloor$ registers inside the block. Average over the uniform block. The expected number of active registers is $qp$, so $\mathbb E\|(I-\Pi)\psi\|^2\le2qp/d$. Replacing $|\psi\rangle$ by the unnormalized $\Pi|\psi\rangle$ changes the acceptance probability by at most $2\|(I-\Pi)\psi\|$, even after the oracle and the final measurement. Once projected, the average output density matrix is identical for the two word laws at this fixed block. Each matrix entry involves one ket query tuple and one bra query tuple, which together consult at most $d$ word coordinates. The two laws have the same joint marginal there. Hence the unprojected acceptance gap is at most $4\sqrt{2qp/d}$. A bounded-error tester on our promised inputs has a constant gap, implying $q=\Omega(d/p)$.

If no such $M$ fits, then $N<8T$. Use the code on all $N$ entries with $\rho$ small enough that every word costs at most $s$. By (1) and the fixed Shannon margin a uniform word is $\xi_0 N$-far with probability $1-o(1)$ when $H_2(\xi_0)<\epsilon/2$. Its marginals and the codeword's agree on every $\delta N$ coordinates, giving $\Omega(N)$ even adaptively by the polynomial method. The constructions use exact members on the easy side, and so the lower bound applies to every stated tolerance.
\end{proof}
In the exact total case a general theorem of Montanaro already forces linear nonadaptive quantum query complexity for every Boolean decision function depending on all $N$ entries \cite{mont}. The hidden block deals with a Hamming-distance promise, where neighboring opposite inputs need not exist. In the range shared with Theorem~\ref{th:testing}, the two results quantify the advantage of adaptive queries at shrinking distance.

\section{Testing with one clean qubit}\label{sec:clean}
Recent work of Shao \cite{shao} studies quantum query algorithms with one initially pure qubit, any number of maximally mixed qubits, input-independent unitaries, controlled phase queries, and a final measurement. We use this model with unbounded mixed workspace. Its queries can be interleaved with arbitrary computation. Limiting the initial purity changes the answer to our testing problem.
\begin{theorem}
In the range of Theorem~\ref{th:nonadaptive}, the bounded-error one-clean-qubit query cost of testing Boolean MCSP is
\[
                  \Theta_{\eta,\epsilon,\theta}\bigl(\min\{N,T/\xi\}\bigr).
\]
The result holds for the exact yes side and for the tolerant yes side specified there.
\end{theorem}

\paragraph{A lower bound for inputs supported on a block.}
Let $g$ be a partial Boolean function on $M$ bits. In an oracle of length $N$, suppose a known set $B$ contains copies of the $M$ bits according to a known map and all positions outside $B$ are fixed. Put $p=|B|/N$. If $g$ has ordinary bounded-error quantum query complexity $d$, evaluating this embedded problem with $k\ge1$ initially pure qubits and all remaining qubits maximally mixed takes
\[
                    \Omega\!\left(\frac{d}{\min\{1,2^k p\}}\right)
                    \quad\text{queries}.                                      \tag{3}
\]
The final measurement can be arbitrary. To see the bound, use the general adversary characterization for partial Boolean functions \cite{adv}. Take a Hermitian adversary matrix $\Gamma$, zero on equal output labels, and normalize $\max_j\|\Gamma\circ\Delta_j\|=1$, where $\Delta_j(x,y)=\mathbf1[x_j\ne y_j]$. We can arrange $\|\Gamma\|=\Omega(d)$ and take a unit eigenvector $v$ for a positive largest eigenvalue, changing the sign if needed.

Purify the same mixed initial state for every input. The ordinary adversary progress is the average Gram inner product of the resulting states weighted by $\overline v_x\Gamma_{xy}v_y$. Its initial value is $\|\Gamma\|$. A successful final measurement changes it by $\Omega(\|\Gamma\|)$ by the usual bounded-error adversary argument \cite{adv}. Write $\psi_{x,i}$ for the branch of a prequery state at address $i$ with phase control active. A controlled phase query changes progress in absolute value by at most
\[
       2\sum_{i\in B}\sum_x |v_x|^2\|\psi_{x,i}\|^2 .
\]
Indeed, at each $i$ the signed quadratic form has matrix $\Gamma\circ\Delta_{j(i)}$ tensored with the identity on the workspace, of norm at most one. Fixed addresses make zero contribution. On a computational space of dimension $D$, the state on any particular input still has operator norm $2^k/D$ after all preceding unitaries. The projector onto addresses in $B$ has rank $pD$, even if we ignore the phase-control condition. Its mass is at most $\min(1,2^kp)$ for each $x$. Telescoping proves (3).

Use the hard base promise from Theorem~\ref{th:nonadaptive}: a codeword versus a word $\alpha M$-far from every size-$s$ table. Its ordinary $d=\Omega(M)$ follows because the easy law has $\Omega(M)$-wise uniform bits and a uniformly random word is far with probability $1-o(1)$. No polynomial of smaller degree can distinguish these unconditional laws by a constant. For (3) the active prefix may even be fixed in advance. Taking $k=1$ gives $\Omega(M/p)$, as required. The dense case gives $\Omega(N)$ already in the ordinary model.

\paragraph{The matching upper bound.}
Shao's Lemma~3.8 computes an arbitrary Boolean predicate of $q$ bits sampled at independent uniform addresses with $2q$ queries and one clean qubit \cite{shao}. The postprocessing predicate can also depend on the addresses. In the proof, all address registers are preserved on each computational branch, so the input-independent reversible postprocessing can select its predicate according to that branch. This implements the sample tester in Theorem~\ref{th:nonadaptive}. To read the entire table, start with one maximally mixed address $J$. For each $i\in\{0,1\}^n$, temporarily replace $J$ by $J\oplus i$, query it into the corresponding mixed mask bit, and undo the address change. Perform this loading and unloading in the same two places in Shao's construction. Conditional on $J$, the $N$ loaded bits give the whole input in a known order. Use the acceptance predicate in that order. This gives $2N$ queries. The workspace and intervening computation may be very large.

\section{Previous work and scope}
Minsky--Papert symmetrization and polynomial methods for quantum query lower bounds motivate the scalar interpolation. High-order polynomial growth for search is classical \cite{ksw}. Kothari supplies the upper bound for identification, and Ambainis and coauthors studied worst-case small on-sets of a given cardinality \cite{small}. Earlier MCSP lower bounds use local pseudorandom generators \cite{cklm}. The additional step here is to compare two shifted copies of a sparse mask. A row can encode a full word of shift information while every monomial that depends on a new row introduces at least one power of the thinning parameter. The theorem concerns the ordinary approximate degree of the classical total decision function. A polynomial that is required to approximate only sparse inputs can take arbitrarily large values elsewhere. The bound on the full Boolean cube in Lemma~\ref{lem:analytic} is essential. A companion argument gives quantum query complexity and approximate degree $\Theta(\sqrt N)$ through $s\le(1+c)n$ for one positive absolute $c$, for either gate basis. The transition between this result and the present theorem remains open here.
\begin{thebibliography}{9}
\bibitem{dw} R. de Wolf. \href{https://arxiv.org/abs/0802.1816}{A note on quantum algorithms and the minimal degree of $\epsilon$-error polynomials for symmetric functions}. Quantum Inf. Comput. 8 (2008), 943--950.
\bibitem{shao} C. Shao. \href{https://arxiv.org/abs/2611.26231}{A sensitivity-constrained polynomial method for DQC1 query complexity}. Preprint, 19 November 2026.
\bibitem{adv} T. Lee, R. Mittal, B. W. Reichardt, R. \v{S}palek and M. Szegedy. \href{https://arxiv.org/abs/1011.3020}{Quantum query complexity of state conversion}. FOCS, 2011.
\bibitem{mont} A. Montanaro. \href{https://arxiv.org/abs/1001.0018}{Nonadaptive quantum query complexity}. Inform. Process. Lett. 110 (2010), 1110--1113.
\bibitem{koth} R. Kothari. \href{https://doi.org/10.4230/LIPIcs.STACS.2014.482}{An optimal quantum algorithm for the oracle identification problem}. STACS, 2014.
\bibitem{beals} R. Beals, H. Buhrman, R. Cleve, M. Mosca and R. de Wolf. \href{https://arxiv.org/abs/quant-ph/9802049}{Quantum lower bounds by polynomials}. J. ACM 48 (2001).
\bibitem{ksw} H. Klauck, R. \v{S}palek and R. de Wolf. \href{https://homepages.cwi.nl/~rdewolf/publ/qc/ksw-sicomp.pdf}{Quantum and classical strong direct product theorems and optimal time-space tradeoffs}. SIAM J. Comput. 36 (2007).
\bibitem{lup} O. B. Lupanov. On a certain approach to the synthesis of control systems: the principle of local coding. Problemy Kibernetiki 14 (1965), 31--110.
\bibitem{kors} A. D. Korshunov. \href{https://www.mathnet.ru/eng/rm9459}{Computational complexity of Boolean functions}. Russian Math. Surveys 67 (2012), Section 2.4.
\bibitem{chash} A. V. Chashkin. \href{https://keldysh.ru/papers/2024/mvk/mvk2024_152.pdf}{Mathematical Questions of Cybernetics (2024)}, Theorem 6.1.
\bibitem{small} A. Ambainis et al. \href{https://www.kecl.ntt.co.jp/people/tani.seiichiro/pdf/SmallOnSets.pdf}{Quantum query complexity of Boolean functions with small on-sets}. ISAAC, 2008.
\bibitem{cklm} M. Cheraghchi, V. Kabanets, Z. Lu and D. Myrisiotis. \href{https://doi.org/10.4230/LIPIcs.ICALP.2019.39}{Circuit lower bounds for MCSP from local pseudorandom generators}. ICALP, 2019.
\bibitem{iso} N. Alon, E. Blais, S. Chakraborty, D. Garc\'ia-Soriano and A. Matsliah. \href{https://doi.org/10.1137/110832677}{Nearly tight bounds for testing function isomorphism}. SIAM J. Comput., 2013.
\bibitem{extractors} S. Chakraborty, D. Garc\'ia-Soriano and A. Matsliah. \href{https://www.isical.ac.in/~sourav/papers/implicit.pdf}{Efficient sample extractors for juntas with applications}.
\end{thebibliography}
\end{document}
